% !TEX root = ../../../main.tex
\section{Continuity}
\label{sec:analysis-i:continuity}

% Source: images/LEC05/IMG_9250.jpg and IMG_9251.jpg.
% Composition is stated in images/LEC05/IMG_9258.jpg.
% The source only sketches the irrational-continuity argument for Thomae's
% function. Added justifications and prerequisites are marked for review.

\subsection{Limits and continuity at a point}

Let $I\subseteq\R$ be an interval
\pendingreviewtext{containing at least two points}, and let $f:I\to\R$.

\begin{definition}[Limit of a function]
  \label{thm:analysis-i:function-limit}
  \pendingreviewtext{Suppose that $c$ can be approached by points of
  $I\setminus\{c\}$.} We write
  \[
    \lim_{x\to c}f(x)=L
  \]
  if, for every $\varepsilon>0$, there is a $\delta>0$ such that
  \[
    x\in I,\quad 0<|x-c|<\delta
    \quad\Longrightarrow\quad |f(x)-L|<\varepsilon.
  \]
\end{definition}

\begin{pendingreview}
The condition $0<|x-c|$ excludes the point $c$ itself. The limit describes
the values near $c$, so it does not depend on the value $f(c)$, or even
require that $f(c)$ be defined.
\end{pendingreview}

The \emph{left-hand limit} and \emph{right-hand limit} are written
\[
  \lim_{x\to c^-}f(x)=L_-,
  \qquad
  \lim_{x\to c^+}f(x)=L_+.
\]
\pendingreviewtext{Their definitions restrict the input condition to
$0<c-x<\delta$ and $0<x-c<\delta$, respectively. When $c$ is an interior
point of $I$, the two-sided limit exists exactly when both one-sided
limits exist and have the same value.}

\begin{definition}[Continuity]
  \label{thm:analysis-i:continuity-at-a-point}
  The function $f$ is \emph{continuous at} $c\in I$ if
  \[
    \lim_{x\to c}f(x)=f(c).
  \]
  \pendingreviewtext{Equivalently, for every $\varepsilon>0$, there is a
  $\delta>0$ such that}
  \[
    \pendingreviewtext{x\in I,\quad |x-c|<\delta
    \quad\Longrightarrow\quad |f(x)-f(c)|<\varepsilon.}
  \]
  It is \emph{continuous on} $I$ if it is continuous at every point of $I$.
\end{definition}

\begin{pendingreview}
For $c\in I$, \emph{left continuity} and \emph{right continuity}
mean, respectively,
\[
  \lim_{x\to c^-}f(x)=f(c),\qquad
  \lim_{x\to c^+}f(x)=f(c).
\]
Each condition is considered when $c$ can be approached from that
side through the domain.
At an endpoint of $I$, only inputs belonging to $I$ are considered.
Thus continuity on $[a,b]$ means ordinary continuity in $(a,b)$,
right continuity at $a$, and left continuity at $b$.
\autoref{fig:analysis-i:continuity-basic-neighbourhood} shows how an input
neighbourhood is chosen to fit a prescribed output neighbourhood.
\end{pendingreview}

\begin{marginfigure}
  \centering
  % Schematic continuity graph; all strokes inherit the shared 0.5pt weight.
\begin{tikzpicture}[analysis diagram,x=9mm,y=9mm]
  \fill[gray!10] (1.25,0) rectangle (2.35,2.65);
  \fill[gray!15] (0,0.9684) rectangle (3.45,2.1684);
  \draw[analysis guide] (0,0.9684) -- (3.45,0.9684);
  \draw[analysis guide] (0,2.1684) -- (3.45,2.1684);
  \draw[analysis guide] (1.25,0) -- (1.25,2.65);
  \draw[analysis guide] (2.35,0) -- (2.35,2.65);
  \draw[->] (-0.12,0) -- (3.7,0) node[right=4pt] {$x$};
  \draw[->] (0,-0.08) -- (0,3.35) node[above=4pt] {$y$};
  \draw[<->,domain=0.12:3.18,samples=45]
    plot (\x,{0.6+0.25*\x+0.16*\x*\x});
  \draw[gray,densely dashed] (0,1.5684) -- (1.8,1.5684) -- (1.8,0);
  \fill (1.8,1.5684) circle[radius=1.6pt];
  \node[left=5pt] at (0,1.5684) {$f(c)$};
  \node[below=5pt] at (1.8,0) {$c$};
  \draw[<->,gray] (3.45,0.9684) -- (3.45,2.1684)
    node[midway,right=4pt,text=black] {$2\varepsilon$};
  \draw[<->,gray] (1.25,-0.55) -- (2.35,-0.55)
    node[midway,below=4pt,text=black] {$2\delta$};
\end{tikzpicture}

  \caption[Input and output neighbourhoods for continuity.]{%
    \pendingreviewtext{A schematic continuous graph. Inputs within
    $\delta$ of $c$ have outputs within $\varepsilon$ of $f(c)$.
    The graph over the shaded input strip lies inside the horizontal
    output band.}}
  \label{fig:analysis-i:continuity-basic-neighbourhood}
\end{marginfigure}

\subsection{Algebra of limits}

\begin{proposition}[Limit laws]
  \label{thm:analysis-i:function-limit-laws}
  Suppose that
  \[
    \lim_{x\to c}f(x)=F,
    \qquad
    \lim_{x\to c}g(x)=G,
  \]
  where $F,G\in\R$. For every $k\in\R$,
  \begin{align*}
    \lim_{x\to c}k f(x)&=kF,\\
    \lim_{x\to c}\bigl(f(x)+g(x)\bigr)&=F+G,\\
    \lim_{x\to c}\bigl(f(x)-g(x)\bigr)&=F-G,\\
    \lim_{x\to c}f(x)g(x)&=FG.
  \end{align*}
  Also,
  \[
    \lim_{x\to c}\frac{f(x)}{g(x)}=\frac{F}{G}
    \qquad\pendingreviewtext{\text{provided }}G\ne0.
  \]
\end{proposition}
% IMG_9251 writes "iff lim g(x) != 0". The nonzero condition is sufficient
% for this quotient law, not necessary for a quotient to have a limit.

\begin{pendingreview}
The last condition ensures that $g(x)$ is nonzero sufficiently near $c$.
These laws show that sums, scalar multiples, and products of continuous
functions are continuous. A quotient is continuous wherever its denominator
is nonzero. In particular, constant functions and $x\mapsto x$ are
continuous, so every polynomial is continuous.
\end{pendingreview}

% !TEX root = ../../hw05.tex
% Source: MA-GY-6213-hw5.pdf, page 1, question 6.
\begin{exercise}
  \label{ex:hw05:absolute-value-continuity}
  Use the definition of a function being continuous at a point, show
  that if $f(x)$ is continuous at $x_0$, so is the function
  \[
    g(x)=|f(x)|.
  \]
  \solutionlink{hw05:absolute-value-continuity}
\end{exercise}


\subsection{Testing continuity with sequences}

\begin{pendingreview}
The following characterization connects the definition above with the
sequence methods already developed. The sequences approach $c$ through
the domain of the function, including from only one side at an endpoint.

\begin{theorem}[Sequential characterization of continuity]
  \label{thm:analysis-i:sequential-continuity}
  Let $f:I\to\R$ and $c\in I$. Then $f$ is continuous at $c$ if and only if
  every sequence $(x_n)$ in $I$ satisfies
  \[
    x_n\to c\quad\Longrightarrow\quad f(x_n)\to f(c).
  \]
\end{theorem}

\begin{proof}
  Suppose first that $f$ is continuous at $c$ and that $x_n\to c$.
  Given $\varepsilon>0$, continuity provides a $\delta>0$ such that
  $|f(x)-f(c)|<\varepsilon$ whenever $x\in I$ and $|x-c|<\delta$.
  Since $x_n\to c$, this input condition holds for all sufficiently large
  $n$. Hence $f(x_n)\to f(c)$.

  Conversely, suppose that $f$ is not continuous at $c$. There is then
  an $\varepsilon_0>0$ such that, for every $n\in\N$, we can choose
  $x_n\in I$ with
  \[
    |x_n-c|<\frac1n,
    \qquad |f(x_n)-f(c)|\ge\varepsilon_0.
  \]
  Thus $x_n\to c$, while $f(x_n)$ does not tend to $f(c)$. This contradicts
  the stated sequence property and proves the converse.
\end{proof}
\end{pendingreview}

% !TEX root = ../../hw05.tex
% Source: MA-GY-6213-hw5.pdf, page 1, question 4.
\begin{exercise}
  \label{ex:hw05:continuous-additive-function}
  \solutionlink{hw05:continuous-additive-function}
  Let $f(x)$ be continuous over $\R$ satisfying the additivity condition
  \[
    f(s+t)=f(s)+f(t),\qquad s,t\in\R.
  \]
  Show that there is a constant $a$ such that $f(x)=ax$.
  (Hint: First prove this when $x$ is an integer and identify what $a$
  is. Then prove this when $x$ is rational. Finally pass to limit for
  any real $x$ by using the denseness of $\Q$ in $\R$.)
\end{exercise}


\subsection{Examples of discontinuity}

\begin{definition}[Indicator function]
  \label{thm:analysis-i:indicator-function}
  For a set $S\subseteq\R$, its \emph{indicator function} is
  \[
    \mathbf{1}_S(x)=
    \begin{cases}
      0,&x\notin S,\\
      1,&x\in S.
    \end{cases}
  \]
\end{definition}

\begin{example}[A jump at one point]
  \label{thm:analysis-i:step-discontinuity}
  Fix $a\in\R$ and define
  \[
    H(x)=
    \begin{cases}
      0,&x\le a,\\
      1,&x>a.
    \end{cases}
  \]
  \pendingreviewtext{Its one-sided limits at $a$ are different, since}
  \[
    \pendingreviewtext{\lim_{x\to a^-}H(x)=0,
    \qquad \lim_{x\to a^+}H(x)=1.}
  \]
  \pendingreviewtext{Consequently $H$ is discontinuous at $a$, although
  it is constant near every other point and hence continuous there.}
\end{example}

\begin{example}[The Dirichlet function]
  \label{thm:analysis-i:dirichlet-function}
  The \emph{Dirichlet function} is
  \[
    D(x)=\mathbf{1}_{\Q}(x)=
    \begin{cases}
      0,&x\notin\Q,\\
      1,&x\in\Q.
    \end{cases}
  \]
  This function is nowhere continuous.

  \begin{pendingreview}
  Every neighbourhood of a real number contains both rational and
  irrational numbers. Thus, however small the neighbourhood, the function
  takes both values $0$ and $1$. Taking $\varepsilon=1/2$ in the definition
  rules out continuity at every point.
  \end{pendingreview}
\end{example}

\begin{example}[Thomae's function]
  \label{thm:analysis-i:thomae-function}
  \emph{Thomae's function}, also called the \emph{popcorn function}, is
  \[
    T(x)=
    \begin{cases}
      0,&x\notin\Q,\\
      \dfrac1q,&x=\dfrac pq\in\Q.
    \end{cases}
  \]
  \pendingreviewtext{In the rational case, the fraction is written in
  lowest terms, with $p\in\Z$, $q\in\N$, and $q>0$.}
  In particular, $T(0)=1$.

  The function $T$ is discontinuous at every rational point.
  \pendingreviewtext{Indeed, if $c=p/q$ is rational in lowest terms,
  then $T(c)=1/q>0$. Choose irrational numbers $x_n\to c$.
  Since $T(x_n)=0$ for every $n$, the sequential characterization
  rules out continuity at $c$.}

  The function $T$ is continuous at every irrational point $c$.
  For a sequence of rational numbers $r_k=p_k/q_k\to c$, the denominators
  tend to infinity, giving
  \[
    T(r_k)=\frac1{q_k}\longrightarrow0=T(c).
  \]

  \begin{pendingreview}
  To justify the assertion about denominators, fix $M\in\N$.
  The convergent sequence $(r_k)$ eventually lies in $[c-1,c+1]$.
  Only finitely many reduced fractions in this interval have denominator
  at most $M$, since their numerators are also bounded. None equals the
  irrational number $c$, so a sufficiently small neighbourhood of $c$
  avoids all of them. Eventually $q_k>M$, proving that $q_k\to\infty$.

  Now take any sequence $x_n\to c$. Its irrational terms have
  $T(x_n)=0$, while the argument above makes the values at its rational
  terms arbitrarily small as well. Hence $T(x_n)\to0=T(c)$, and
  \autoref{thm:analysis-i:sequential-continuity} proves continuity at $c$.
  \autoref{fig:analysis-i:continuity-basic-thomae} illustrates why only
  finitely many large values need to be avoided near $c$.
  \end{pendingreview}
\end{example}

\begin{marginfigure}
  \centering
  % Only selected rational values are drawn, not a connected approximation.
\begin{tikzpicture}[analysis diagram,x=35mm,y=22mm]
  \draw[->] (-0.035,0) -- (1.13,0) node[right=4pt] {$x$};
  \draw[->] (0,-0.04) -- (0,1.17) node[above=4pt] {$T(x)$};
  \draw[gray] (-0.02,0.5) -- (0.02,0.5);
  \node[left=4pt] at (0,0.5) {$\frac12$};
  \node[left=4pt] at (0,1) {$1$};
  \node[below=5pt] at (0,0) {$0$};
  \draw (1,-0.018) -- (1,0.018);
  \node[below=5pt] at (1,0) {$1$};
  \foreach \p/\q in {
    0/1,1/1,1/2,1/3,2/3,1/4,3/4,1/5,2/5,3/5,4/5,
    1/6,5/6,1/7,2/7,3/7,4/7,5/7,6/7,1/8,3/8,5/8,7/8
  } {
    \fill ({\p/\q},{1/\q}) circle[radius=1.2pt];
  }
\end{tikzpicture}

  \caption[Selected values of Thomae's function.]{%
    \pendingreviewtext{Selected values of $T$ on $[0,1]$. Dots show all
    reduced fractions with denominator at most $8$. Larger denominators
    give smaller heights; irrational inputs have value $0$. The finite
    drawing shows only part of the graph.}}
  \label{fig:analysis-i:continuity-basic-thomae}
\end{marginfigure}

\subsection{Composition}

% Source: images/LEC05/IMG_9258.jpg, paragraph headed "Composition".
\begin{theorem}[Composition of continuous functions]
  \label{thm:analysis-i:continuous-composition}
  \pendingreviewtext{Let $f:I\to\R$ and $g:J\to\R$, where $f(I)\subseteq J$.}
  If $f$ is continuous at $c$ and $g$ is continuous at $f(c)$, then
  $g\circ f$ is continuous at $c$, where
  \[
    (g\circ f)(x)=g(f(x)).
  \]
\end{theorem}

\begin{pendingreview}
\begin{proof}
  If $x_n\to c$ with $x_n\in I$, continuity of $f$ gives
  $f(x_n)\to f(c)$. These values lie in $J$, so continuity of $g$ gives
  $g(f(x_n))\to g(f(c))$. Apply
  \autoref{thm:analysis-i:sequential-continuity} to $g\circ f$.
\end{proof}
\end{pendingreview}

% !TEX root = ../../hw05.tex
% Source: MA-GY-6213-hw5.pdf, page 1, question 5.
\begin{exercise}
  \label{ex:hw05:composition-continuity}
  \solutionlink{hw05:composition-continuity}
  Use the definition of a function being continuous at a point, prove
  that if $f(x)$ is continuous at $x_0$, $g(y)$ is continuous at
  $y_0=f(x_0)$, then $h(x)=g(f(x))=(g\circ f)(x)$ is continuous at
  $x_0$ as well.
\end{exercise}

